\documentclass[12pt]{article}
\usepackage[utf8]{inputenc}
\usepackage[T1]{fontenc}
\usepackage{hyperref}
\title{Paunaka Language Resources for K-12 Curriculum Development}
\author{Compiled by Claude Sonnet 5.5 (claude-sonnet-5-5) for Andr\'as Kornai}
\date{}
\begin{document}
\maketitle

\section{Introduction}
Paunaka (ISO 639-3 pnk) is an Arawakan language of the Chiquitania lowlands of Bolivia. Speaker counts in the sources differ: eight in 2023 according to Terhart's grammar, 11 according to the ELDP archive and a 2016 Last Whispers page, fewer than ten in a 2018 paper. The Paunaka Documentation Project (ELDP, 2011--2013) produced about 120 hours of recordings.

\section{Primary Educational Resources}
No primers, textbooks or school programmes for Paunaka were found in the searches.

\section{Linguistic and Reference Materials}
\begin{itemize}
\item Lena Terhart, \emph{A grammar of Paunaka} (Language Science Press, 2024; CC BY 4.0, over 1500 examples): \url{https://langsci-press.org/catalog/book/400}
\item Open-access copy at OAPEN: \url{https://library.oapen.org/handle/20.500.12657/90227?show=full}
\item Repository record at Freie Universit\"at Berlin: \url{https://refubium.fu-berlin.de/handle/fub188/42745?show=full}
\item Journal article on Paunaka (Redalyc; the search did not show its title): \url{https://www.redalyc.org/journal/3940/394056633008/}
\end{itemize}

\section{Religious and Community Resources}
No scripture translations or community organisation pages were found in the searches.

\section{Digital and Audio Resources}
\begin{itemize}
\item Paunaka language archive (ELAR/ELDP), audio and video recordings deposited mainly by Swintha Danielsen: \url{https://eldp.access.preservica.com/dk0104/}
\item Last Whispers page on Paunaka: \url{https://lastwhispers.org/paunaka}
\end{itemize}

\section{Curriculum Development Context}
The language is used by very few elderly speakers and is not known to be taught in schools. The openly licensed grammar and the archived recordings are the main foundation for any revitalisation material; the archive may have access conditions that were not checked.

\end{document}